\begin{table*}[t]
\centering
\caption{Condiciones para la comparación del desempeño entre \(B\) y \(A'\)}
\label{tab:sintesis_desempeno_esquemas}
\small
\renewcommand{\arraystretch}{1.25}
\setlength{\tabcolsep}{6pt}

\begin{tabular}{
@{}
>{\centering\arraybackslash}p{0.21\textwidth}
>{\centering\arraybackslash}p{0.17\textwidth}
>{\raggedright\arraybackslash}p{0.54\textwidth}
@{}}
\toprule
\textbf{Caso}
&
\textbf{Resultado}
&
\textbf{Condiciones}
\\
\hline

\(d_Q>0\),
\(ICG^A<\tau\)
&
\(Q^B\geq Q^{A'}\)
&
Condición~\eqref{eq:condicion_suficiente_incentivo_productivo}, igualdad local de las intensidades sancionadoras, mantenimiento del ordenamiento de incentivos en el tramo relevante y óptimo único.
\\[5pt]

\(d_Q>0\),
\(ICG^A>\tau\)
&
\(Q^B\geq Q^{A'}\)
&
Condición~\eqref{eq:condicion_suficiente_incentivo_productivo}, mantenimiento del ordenamiento de incentivos en el tramo relevante y óptimo único. No se requiere igualar las intensidades sancionadoras porque la multa de \(A'\) se encuentra inactiva.
\\[5pt]

\(d_Q=0\),
\(ICG^A<\tau\)
&
\(Q^B=Q^{A'}\)
&
Igualdad local de las intensidades sancionadoras y óptimo único.
\\[5pt]

\(ICG^A=\tau\), o
\(d_Q=0\) e \(ICG^A>\tau\)
&
\(Q^B\geq Q^{A'}\)
&
La fiscalización separada debe inducir \(I_r^B\geq I_r^{A'}\) para toda \(r\in R^{Q}\), pues la comparación marginal no determina por sí sola el ordenamiento.
\\
\hline
\end{tabular}
\end{table*}